% =============================================================
%  ch26_simulierbarkeit.tex  --  Simulierbarkeit in Gazebo und RViz
% =============================================================

\chapter{Simulierbarkeit in Gazebo und RViz}

Lassen sich Anfahren der Zielposition, Schlagbewegung und die induzierten
Kräfte mit Gazebo/RViz abbilden?

\section{Was Gazebo kann -- und was nicht}
\begin{longtable}{@{}p{7.4cm}p{7.2cm}@{}}
\toprule
\textbf{Gazebo bildet ab (Starrkörperdynamik)} & \textbf{Gazebo bildet \emph{nicht} ab} \\
\midrule
Anfahren der Zielposition (prismatische + rotatorische Gelenke) & Strukturelle Biegung/Durchbiegung der Stäbe \\
Schlagbewegung des Arms      & Riemendehnung als kontinuierliche Elastizität \\
Gelenkmomente/-kräfte (auslesbar) & Modale Resonanz/Eigenschwingung der Struktur \\
Dynamische Kopplung (Armträgheit reagiert auf Portal) & Rahmenflattern, Mikrovibrationen \\
Ball-Schläger-Kontakt + Reaktionskräfte & exakte Spin-/Absprungphysik (nur genähert) \\
Motormoment-Budgetierung      & \\
\bottomrule
\end{longtable}

\begin{warnbox}[title=Der Kernpunkt]
Gazebo behandelt Links als \textbf{starre Körper}. Die Schwingung, die dich
sorgt, kommt aber großteils aus \emph{struktureller Nachgiebigkeit}
(Stabbiegung, Riemendehnung, Rahmenflex) -- starre Links biegen sich nicht,
also zeigt Standard-Gazebo diese Resonanz \emph{nicht}.
\end{warnbox}

\section{Wege, die Nachgiebigkeit doch abzubilden}
\begin{itemize}
\item \textbf{Klumpen-Nachgiebigkeit in Gazebo:} ein steifes Feder-Dämpfer-Gelenk
  zwischen Segmenten einfügen (Riemendehnung als prismatische Feder, oder ein
  nachgiebiges Drehgelenk am Auslegerfuß). Das ist eine \emph{Näherung} der
  ersten Mode, kein vollständiges Modalverhalten.
\item \textbf{Echte Strukturvibration:} FEM / flexible Mehrkörpersimulation --
  SolidWorks Simulation (modal), MATLAB \textbf{Simscape Multibody} mit
  flexiblen Körpern (reduzierte FEM-Modelle), Adams. Co-Simulation mit ROS
  möglich.
\item \textbf{RViz:} reine Visualisierung -- zeigt Bewegung und TF, rechnet
  \emph{keine} Physik. Kräfte kannst du als Marker visualisieren, wenn ein
  anderer Knoten sie berechnet/publiziert.
\end{itemize}

\begin{infobox}[title=Empfohlene Aufteilung]
\textbf{Gazebo} für Regelung, Trajektorientracking, Momentenbudget und
Ball-Schläger-Kontakt. \textbf{FEM/Simscape} für die strukturelle Resonanz
($f_1$, Modenformen). Beides verbindet sich über die Input-Shaping-Frequenz:
FEM liefert $f_1$, der ROS-Regler nutzt es zur Schwingungsunterdrückung.
\end{infobox}
